% ============================================================
%  ECON306: W07D2 — Part I Review and Synthesis; Practice Questions
%  Date: Thursday 27 August 2026
%  NOTE: NEW content — no source frames in old master deck.
%        Built from Course Map LOs (W07D2) and Part I content
%        across W01–W07D1 day files.
% ============================================================

\documentclass[aspectratio=169, 12pt]{beamer}

% ============================================================
%  ECON306: Shared Preamble — The Economics of Health and Education
%  University of Otago — Semester 2, 2026
%
%  SHARED BY ALL ECON306 DAY DECKS.
%  Input this file immediately after \documentclass in each deck:
%    % ============================================================
%  ECON306: Shared Preamble — The Economics of Health and Education
%  University of Otago — Semester 2, 2026
%
%  SHARED BY ALL ECON306 DAY DECKS.
%  Input this file immediately after \documentclass in each deck:
%    % ============================================================
%  ECON306: Shared Preamble — The Economics of Health and Education
%  University of Otago — Semester 2, 2026
%
%  SHARED BY ALL ECON306 DAY DECKS.
%  Input this file immediately after \documentclass in each deck:
%    \input{../Preambles/econ306_preamble.tex}
%  Do NOT include \documentclass or per-deck metadata here.
% ============================================================

% ─── Theme & Colours ────────────────────────────────────────
\usetheme{default}
\usecolortheme{default}
\usefonttheme{professionalfonts}
\definecolor{OtagoBlue}{HTML}{003057}
\definecolor{OtagoGold}{HTML}{A8923A}
\definecolor{TextGrey}{HTML}{3A3A3A}
\definecolor{AccentBlue}{HTML}{2B6CB0}
\definecolor{LightBg}{HTML}{F7F9FC}

\setbeamercolor{frametitle}{fg=OtagoBlue, bg=white}
\setbeamercolor{title}{fg=OtagoBlue}
\setbeamercolor{subtitle}{fg=OtagoGold}
\setbeamercolor{author}{fg=TextGrey}
\setbeamercolor{date}{fg=TextGrey}
\setbeamercolor{institute}{fg=TextGrey}
\setbeamercolor{normal text}{fg=TextGrey, bg=white}
\setbeamercolor{item}{fg=OtagoBlue}
\setbeamercolor{subitem}{fg=AccentBlue}
\setbeamercolor{block title}{bg=OtagoBlue!12, fg=OtagoBlue}
\setbeamercolor{block body}{bg=OtagoBlue!5, fg=TextGrey}
\setbeamercolor{alertblock title}{bg=OtagoGold!20, fg=OtagoBlue}
\setbeamercolor{alertblock body}{bg=OtagoGold!8, fg=TextGrey}
\setbeamercolor{alerted text}{fg=OtagoGold!80!black}
\setbeamercolor{structure}{fg=OtagoBlue}

\setbeamertemplate{frametitle}{%
	\vspace{0.35em}%
	\textbf{\insertframetitle}\par%
	\vspace{0.05em}%
	\textcolor{OtagoGold}{\rule{\textwidth}{0.6pt}}}

\setbeamertemplate{navigation symbols}{}
\setbeamertemplate{footline}{%
	\leavevmode\hbox{%
		\begin{beamercolorbox}[wd=.40\paperwidth,ht=2.5ex,dp=1.5ex,leftskip=.3cm]{author in head/foot}%
			\textcolor{TextGrey}{\footnotesize ECON306 \textbullet\ University of Otago}%
		\end{beamercolorbox}%
		\begin{beamercolorbox}[wd=.47\paperwidth,ht=2.5ex,dp=1.5ex,center]{title in head/foot}%
			\textcolor{TextGrey}{\footnotesize \insertshorttitle}%
		\end{beamercolorbox}%
		\begin{beamercolorbox}[wd=.13\paperwidth,ht=2.5ex,dp=1.5ex,rightskip=.3cm,right]{date in head/foot}%
			\textcolor{TextGrey}{\footnotesize \insertframenumber\,/\,\inserttotalframenumber}%
	\end{beamercolorbox}}\vskip0pt}

\setbeamertemplate{itemize item}{\textcolor{OtagoGold}{\small$\blacktriangleright$}}
\setbeamertemplate{itemize subitem}{\textcolor{AccentBlue}{\small$\triangleright$}}
\setbeamertemplate{enumerate item}{\textcolor{OtagoBlue}{\small\bfseries\insertenumlabel.}}

\setbeamertemplate{title page}{%
	\begin{tikzpicture}[remember picture,overlay]
		\fill[OtagoBlue] (current page.north west)
		rectangle ([yshift=-0.38\paperheight]current page.north east);
	\end{tikzpicture}
	\vspace{0.06\paperheight}
	\begin{beamercolorbox}[wd=\textwidth, sep=0pt]{title}
		{\color{white}\Large\textbf{\inserttitle}}\\[0.4em]
		{\color{OtagoGold}\large\insertsubtitle}
	\end{beamercolorbox}
	\vspace{1.4em}
	\begin{beamercolorbox}[wd=\textwidth, sep=0pt]{author}
		{\color{TextGrey}\normalsize\insertauthor}\\[0.2em]
		{\color{TextGrey}\small\insertinstitute}\\[0.2em]
		{\color{TextGrey}\small\insertdate}
\end{beamercolorbox}}

% ─── Packages ──────────────────────────────────────────────
\usepackage{tikz}
\usepackage{booktabs}
\usepackage{array}
\usepackage{graphicx}
\usepackage{hyperref}
\usepackage{amsmath}
\usepackage{amssymb}
\usepackage{colortbl}
\usepackage{multirow}
\usepackage{xcolor}
\hypersetup{colorlinks=true, linkcolor=AccentBlue, urlcolor=AccentBlue,
	citecolor=AccentBlue, pdfauthor={Austin Henderson},
	pdftitle={ECON306 S2 2026}}

% ─── Images folder ─────────────────────────────────────────
%  Place all slide images in the 'images/' subfolder next to each deck.
%  Usage: \includegraphics[width=0.85\textwidth]{figure_name}
\graphicspath{{images/}}

% ─── Reusable macros ───────────────────────────────────────
% Recap slide at start of each lecture
\newcommand{\recapslide}[1]{%
	\begin{frame}{Recap: What Did We Cover Last Time?}
		\begin{block}{Key points from last lecture}
			\begin{itemize}#1\end{itemize}
		\end{block}
\end{frame}}

% Plan slide at start of each lecture
\newcommand{\planslide}[2]{%
	\begin{frame}{Today's Plan — #1}
		\begin{columns}
			\column{0.6\textwidth}
			\begin{itemize}#2\end{itemize}
			\column{0.38\textwidth}
			\begin{alertblock}{Reading}
				\footnotesize See Brightspace (Aoroa) for required readings before this lecture.
			\end{alertblock}
		\end{columns}
\end{frame}}

% Summary slide at end of each lecture
\newcommand{\summaryside}[1]{%
	\begin{frame}{Lecture Summary}
		\begin{exampleblock}{Key takeaways}
			\begin{itemize}#1\end{itemize}
		\end{exampleblock}
		\vspace{0.5em}
		{\footnotesize\textcolor{OtagoGold}{\textbf{Brightspace (Aoroa):}} slides, readings, and any announcements posted there.}
\end{frame}}

% NZ spotlight box
\newcommand{\nzbox}[2]{%
	\begin{alertblock}{\includegraphics[height=0.8em]{nz_flag_icon}\hspace{0.3em} NZ Spotlight: #1}
		#2
\end{alertblock}}

% Tutorial callout
\newcommand{\tutorialbox}[1]{%
	\begin{block}{\textbf{Tutorial This Week}}
		#1
\end{block}}

% ─── Section dividers ──────────────────────────────────────
\AtBeginSection[]{
	\begin{frame}[plain]
		\begin{tikzpicture}[remember picture,overlay]
			\fill[OtagoBlue!95] (current page.north west)
			rectangle (current page.south east);
			\node[anchor=center, text=white, font=\Large\bfseries,
			text width=0.9\paperwidth, align=center]
			at (current page.center) {\insertsectionhead};
			\node[anchor=south, text=OtagoGold, font=\normalsize,
			text width=0.9\paperwidth, align=center]
			at ([yshift=2.5cm]current page.south) {ECON306 — University of Otago};
		\end{tikzpicture}
\end{frame}}

%  Do NOT include \documentclass or per-deck metadata here.
% ============================================================

% ─── Theme & Colours ────────────────────────────────────────
\usetheme{default}
\usecolortheme{default}
\usefonttheme{professionalfonts}
\definecolor{OtagoBlue}{HTML}{003057}
\definecolor{OtagoGold}{HTML}{A8923A}
\definecolor{TextGrey}{HTML}{3A3A3A}
\definecolor{AccentBlue}{HTML}{2B6CB0}
\definecolor{LightBg}{HTML}{F7F9FC}

\setbeamercolor{frametitle}{fg=OtagoBlue, bg=white}
\setbeamercolor{title}{fg=OtagoBlue}
\setbeamercolor{subtitle}{fg=OtagoGold}
\setbeamercolor{author}{fg=TextGrey}
\setbeamercolor{date}{fg=TextGrey}
\setbeamercolor{institute}{fg=TextGrey}
\setbeamercolor{normal text}{fg=TextGrey, bg=white}
\setbeamercolor{item}{fg=OtagoBlue}
\setbeamercolor{subitem}{fg=AccentBlue}
\setbeamercolor{block title}{bg=OtagoBlue!12, fg=OtagoBlue}
\setbeamercolor{block body}{bg=OtagoBlue!5, fg=TextGrey}
\setbeamercolor{alertblock title}{bg=OtagoGold!20, fg=OtagoBlue}
\setbeamercolor{alertblock body}{bg=OtagoGold!8, fg=TextGrey}
\setbeamercolor{alerted text}{fg=OtagoGold!80!black}
\setbeamercolor{structure}{fg=OtagoBlue}

\setbeamertemplate{frametitle}{%
	\vspace{0.35em}%
	\textbf{\insertframetitle}\par%
	\vspace{0.05em}%
	\textcolor{OtagoGold}{\rule{\textwidth}{0.6pt}}}

\setbeamertemplate{navigation symbols}{}
\setbeamertemplate{footline}{%
	\leavevmode\hbox{%
		\begin{beamercolorbox}[wd=.40\paperwidth,ht=2.5ex,dp=1.5ex,leftskip=.3cm]{author in head/foot}%
			\textcolor{TextGrey}{\footnotesize ECON306 \textbullet\ University of Otago}%
		\end{beamercolorbox}%
		\begin{beamercolorbox}[wd=.47\paperwidth,ht=2.5ex,dp=1.5ex,center]{title in head/foot}%
			\textcolor{TextGrey}{\footnotesize \insertshorttitle}%
		\end{beamercolorbox}%
		\begin{beamercolorbox}[wd=.13\paperwidth,ht=2.5ex,dp=1.5ex,rightskip=.3cm,right]{date in head/foot}%
			\textcolor{TextGrey}{\footnotesize \insertframenumber\,/\,\inserttotalframenumber}%
	\end{beamercolorbox}}\vskip0pt}

\setbeamertemplate{itemize item}{\textcolor{OtagoGold}{\small$\blacktriangleright$}}
\setbeamertemplate{itemize subitem}{\textcolor{AccentBlue}{\small$\triangleright$}}
\setbeamertemplate{enumerate item}{\textcolor{OtagoBlue}{\small\bfseries\insertenumlabel.}}

\setbeamertemplate{title page}{%
	\begin{tikzpicture}[remember picture,overlay]
		\fill[OtagoBlue] (current page.north west)
		rectangle ([yshift=-0.38\paperheight]current page.north east);
	\end{tikzpicture}
	\vspace{0.06\paperheight}
	\begin{beamercolorbox}[wd=\textwidth, sep=0pt]{title}
		{\color{white}\Large\textbf{\inserttitle}}\\[0.4em]
		{\color{OtagoGold}\large\insertsubtitle}
	\end{beamercolorbox}
	\vspace{1.4em}
	\begin{beamercolorbox}[wd=\textwidth, sep=0pt]{author}
		{\color{TextGrey}\normalsize\insertauthor}\\[0.2em]
		{\color{TextGrey}\small\insertinstitute}\\[0.2em]
		{\color{TextGrey}\small\insertdate}
\end{beamercolorbox}}

% ─── Packages ──────────────────────────────────────────────
\usepackage{tikz}
\usepackage{booktabs}
\usepackage{array}
\usepackage{graphicx}
\usepackage{hyperref}
\usepackage{amsmath}
\usepackage{amssymb}
\usepackage{colortbl}
\usepackage{multirow}
\usepackage{xcolor}
\hypersetup{colorlinks=true, linkcolor=AccentBlue, urlcolor=AccentBlue,
	citecolor=AccentBlue, pdfauthor={Austin Henderson},
	pdftitle={ECON306 S2 2026}}

% ─── Images folder ─────────────────────────────────────────
%  Place all slide images in the 'images/' subfolder next to each deck.
%  Usage: \includegraphics[width=0.85\textwidth]{figure_name}
\graphicspath{{images/}}

% ─── Reusable macros ───────────────────────────────────────
% Recap slide at start of each lecture
\newcommand{\recapslide}[1]{%
	\begin{frame}{Recap: What Did We Cover Last Time?}
		\begin{block}{Key points from last lecture}
			\begin{itemize}#1\end{itemize}
		\end{block}
\end{frame}}

% Plan slide at start of each lecture
\newcommand{\planslide}[2]{%
	\begin{frame}{Today's Plan — #1}
		\begin{columns}
			\column{0.6\textwidth}
			\begin{itemize}#2\end{itemize}
			\column{0.38\textwidth}
			\begin{alertblock}{Reading}
				\footnotesize See Brightspace (Aoroa) for required readings before this lecture.
			\end{alertblock}
		\end{columns}
\end{frame}}

% Summary slide at end of each lecture
\newcommand{\summaryside}[1]{%
	\begin{frame}{Lecture Summary}
		\begin{exampleblock}{Key takeaways}
			\begin{itemize}#1\end{itemize}
		\end{exampleblock}
		\vspace{0.5em}
		{\footnotesize\textcolor{OtagoGold}{\textbf{Brightspace (Aoroa):}} slides, readings, and any announcements posted there.}
\end{frame}}

% NZ spotlight box
\newcommand{\nzbox}[2]{%
	\begin{alertblock}{\includegraphics[height=0.8em]{nz_flag_icon}\hspace{0.3em} NZ Spotlight: #1}
		#2
\end{alertblock}}

% Tutorial callout
\newcommand{\tutorialbox}[1]{%
	\begin{block}{\textbf{Tutorial This Week}}
		#1
\end{block}}

% ─── Section dividers ──────────────────────────────────────
\AtBeginSection[]{
	\begin{frame}[plain]
		\begin{tikzpicture}[remember picture,overlay]
			\fill[OtagoBlue!95] (current page.north west)
			rectangle (current page.south east);
			\node[anchor=center, text=white, font=\Large\bfseries,
			text width=0.9\paperwidth, align=center]
			at (current page.center) {\insertsectionhead};
			\node[anchor=south, text=OtagoGold, font=\normalsize,
			text width=0.9\paperwidth, align=center]
			at ([yshift=2.5cm]current page.south) {ECON306 — University of Otago};
		\end{tikzpicture}
\end{frame}}

%  Do NOT include \documentclass or per-deck metadata here.
% ============================================================

% ─── Theme & Colours ────────────────────────────────────────
\usetheme{default}
\usecolortheme{default}
\usefonttheme{professionalfonts}
\definecolor{OtagoBlue}{HTML}{003057}
\definecolor{OtagoGold}{HTML}{A8923A}
\definecolor{TextGrey}{HTML}{3A3A3A}
\definecolor{AccentBlue}{HTML}{2B6CB0}
\definecolor{LightBg}{HTML}{F7F9FC}

\setbeamercolor{frametitle}{fg=OtagoBlue, bg=white}
\setbeamercolor{title}{fg=OtagoBlue}
\setbeamercolor{subtitle}{fg=OtagoGold}
\setbeamercolor{author}{fg=TextGrey}
\setbeamercolor{date}{fg=TextGrey}
\setbeamercolor{institute}{fg=TextGrey}
\setbeamercolor{normal text}{fg=TextGrey, bg=white}
\setbeamercolor{item}{fg=OtagoBlue}
\setbeamercolor{subitem}{fg=AccentBlue}
\setbeamercolor{block title}{bg=OtagoBlue!12, fg=OtagoBlue}
\setbeamercolor{block body}{bg=OtagoBlue!5, fg=TextGrey}
\setbeamercolor{alertblock title}{bg=OtagoGold!20, fg=OtagoBlue}
\setbeamercolor{alertblock body}{bg=OtagoGold!8, fg=TextGrey}
\setbeamercolor{alerted text}{fg=OtagoGold!80!black}
\setbeamercolor{structure}{fg=OtagoBlue}

\setbeamertemplate{frametitle}{%
	\vspace{0.35em}%
	\textbf{\insertframetitle}\par%
	\vspace{0.05em}%
	\textcolor{OtagoGold}{\rule{\textwidth}{0.6pt}}}

\setbeamertemplate{navigation symbols}{}
\setbeamertemplate{footline}{%
	\leavevmode\hbox{%
		\begin{beamercolorbox}[wd=.40\paperwidth,ht=2.5ex,dp=1.5ex,leftskip=.3cm]{author in head/foot}%
			\textcolor{TextGrey}{\footnotesize ECON306 \textbullet\ University of Otago}%
		\end{beamercolorbox}%
		\begin{beamercolorbox}[wd=.47\paperwidth,ht=2.5ex,dp=1.5ex,center]{title in head/foot}%
			\textcolor{TextGrey}{\footnotesize \insertshorttitle}%
		\end{beamercolorbox}%
		\begin{beamercolorbox}[wd=.13\paperwidth,ht=2.5ex,dp=1.5ex,rightskip=.3cm,right]{date in head/foot}%
			\textcolor{TextGrey}{\footnotesize \insertframenumber\,/\,\inserttotalframenumber}%
	\end{beamercolorbox}}\vskip0pt}

\setbeamertemplate{itemize item}{\textcolor{OtagoGold}{\small$\blacktriangleright$}}
\setbeamertemplate{itemize subitem}{\textcolor{AccentBlue}{\small$\triangleright$}}
\setbeamertemplate{enumerate item}{\textcolor{OtagoBlue}{\small\bfseries\insertenumlabel.}}

\setbeamertemplate{title page}{%
	\begin{tikzpicture}[remember picture,overlay]
		\fill[OtagoBlue] (current page.north west)
		rectangle ([yshift=-0.38\paperheight]current page.north east);
	\end{tikzpicture}
	\vspace{0.06\paperheight}
	\begin{beamercolorbox}[wd=\textwidth, sep=0pt]{title}
		{\color{white}\Large\textbf{\inserttitle}}\\[0.4em]
		{\color{OtagoGold}\large\insertsubtitle}
	\end{beamercolorbox}
	\vspace{1.4em}
	\begin{beamercolorbox}[wd=\textwidth, sep=0pt]{author}
		{\color{TextGrey}\normalsize\insertauthor}\\[0.2em]
		{\color{TextGrey}\small\insertinstitute}\\[0.2em]
		{\color{TextGrey}\small\insertdate}
\end{beamercolorbox}}

% ─── Packages ──────────────────────────────────────────────
\usepackage{tikz}
\usepackage{booktabs}
\usepackage{array}
\usepackage{graphicx}
\usepackage{hyperref}
\usepackage{amsmath}
\usepackage{amssymb}
\usepackage{colortbl}
\usepackage{multirow}
\usepackage{xcolor}
\hypersetup{colorlinks=true, linkcolor=AccentBlue, urlcolor=AccentBlue,
	citecolor=AccentBlue, pdfauthor={Austin Henderson},
	pdftitle={ECON306 S2 2026}}

% ─── Images folder ─────────────────────────────────────────
%  Place all slide images in the 'images/' subfolder next to each deck.
%  Usage: \includegraphics[width=0.85\textwidth]{figure_name}
\graphicspath{{images/}}

% ─── Reusable macros ───────────────────────────────────────
% Recap slide at start of each lecture
\newcommand{\recapslide}[1]{%
	\begin{frame}{Recap: What Did We Cover Last Time?}
		\begin{block}{Key points from last lecture}
			\begin{itemize}#1\end{itemize}
		\end{block}
\end{frame}}

% Plan slide at start of each lecture
\newcommand{\planslide}[2]{%
	\begin{frame}{Today's Plan — #1}
		\begin{columns}
			\column{0.6\textwidth}
			\begin{itemize}#2\end{itemize}
			\column{0.38\textwidth}
			\begin{alertblock}{Reading}
				\footnotesize See Brightspace (Aoroa) for required readings before this lecture.
			\end{alertblock}
		\end{columns}
\end{frame}}

% Summary slide at end of each lecture
\newcommand{\summaryside}[1]{%
	\begin{frame}{Lecture Summary}
		\begin{exampleblock}{Key takeaways}
			\begin{itemize}#1\end{itemize}
		\end{exampleblock}
		\vspace{0.5em}
		{\footnotesize\textcolor{OtagoGold}{\textbf{Brightspace (Aoroa):}} slides, readings, and any announcements posted there.}
\end{frame}}

% NZ spotlight box
\newcommand{\nzbox}[2]{%
	\begin{alertblock}{\includegraphics[height=0.8em]{nz_flag_icon}\hspace{0.3em} NZ Spotlight: #1}
		#2
\end{alertblock}}

% Tutorial callout
\newcommand{\tutorialbox}[1]{%
	\begin{block}{\textbf{Tutorial This Week}}
		#1
\end{block}}

% ─── Section dividers ──────────────────────────────────────
\AtBeginSection[]{
	\begin{frame}[plain]
		\begin{tikzpicture}[remember picture,overlay]
			\fill[OtagoBlue!95] (current page.north west)
			rectangle (current page.south east);
			\node[anchor=center, text=white, font=\Large\bfseries,
			text width=0.9\paperwidth, align=center]
			at (current page.center) {\insertsectionhead};
			\node[anchor=south, text=OtagoGold, font=\normalsize,
			text width=0.9\paperwidth, align=center]
			at ([yshift=2.5cm]current page.south) {ECON306 — University of Otago};
		\end{tikzpicture}
\end{frame}}


\title[ECON306: W07D2 — Part I Review]{ECON306\\The Economics of Health and Education}
\subtitle{W07D2: Part I Review and Synthesis; Practice Questions}
\author{Austin Henderson}
\institute{Department of Economics\\University of Otago}
\date{Thursday 27 August 2026}

\begin{document}

\begin{frame}[plain]
  \titlepage
\end{frame}

\recapslide{
  \item NZ health system: 1938 to Te Whatu Ora 2022 --- Beveridge-style, public funding,
        centralised purchasing (PHARMAC)
  \item Each major reform driven by market-failure logic: reduce fragmentation, improve
        incentives, address equity gaps
  \item Current system priorities: access, equity, integration; ongoing fiscal pressure
}

\planslide{Part I Review and Practice Questions}{
  \item Synthesis: the six characteristics, demand models, market failures, insurance,
        providers, systems
  \item Common exam errors and misconceptions to avoid
  \item Practice questions in test format
  \item Tomorrow: \textbf{In-class test (15\%)} --- test-only session, no lecture
}

% ----------------------------------------------------------------
\begin{frame}{What Part I Covered: An Integrated Map}
  \begin{columns}
    \column{0.32\textwidth}
      \begin{block}{Weeks 1--3}
        \footnotesize\textbf{Why is health care special?}
        \begin{itemize}\footnotesize
          \item \textbf{W1}: Six characteristics; government involvement
          \item \textbf{W1--2}: Health production fn; Grossman; Wagstaff
          \item \textbf{W3}: Market failures; SID; regulation theories
        \end{itemize}
      \end{block}
    \column{0.32\textwidth}
      \begin{block}{Weeks 4--5}
        \footnotesize\textbf{Who demands, who funds?}
        \begin{itemize}\footnotesize
          \item \textbf{W4}: Ageing; NZ Super; prevention vs.\ cure
          \item \textbf{W5}: Insurance; moral hazard; adverse selection; NZ hybrid
        \end{itemize}
      \end{block}
    \column{0.32\textwidth}
      \begin{alertblock}{Weeks 6--7}
        \footnotesize\textbf{Providers and systems}
        \begin{itemize}\footnotesize
          \item \textbf{W6}: Hospital behaviour; system types; NZ history; US paradox
          \item \textbf{W7D1}: NZ system depth; structural reasoning
        \end{itemize}
      \end{alertblock}
  \end{columns}
  \vspace{0.3em}
  \begin{exampleblock}{The recurring theme}
    \small The \textbf{six special characteristics} (uncertainty, derived demand, asymmetric information, catastrophic consequences, externalities, price discrimination) underpin every topic in Part I.
  \end{exampleblock}
\end{frame}

% ----------------------------------------------------------------
\begin{frame}{Review: The Six Special Characteristics}
  \small
  \begin{enumerate}
    \item \textbf{Uncertainty} --- demand is unpredictable; insurance is the institutional response
    \item \textbf{Derived demand} --- people want health, not health care per se; $H = f(\text{HC}, \ldots)$
    \item \textbf{Asymmetric information} --- principal-agent problem; SID hypothesis; adverse selection
    \item \textbf{Catastrophic and irreversible consequences} --- outcomes can be life-threatening; risk of catastrophe $\neq$ expected value
    \item \textbf{Externalities} --- physical (immunisation, sanitation) and caring (psychic)
    \item \textbf{Price discrimination} --- insurers use third-degree discrimination; regulation governs this
  \end{enumerate}
  \vspace{0.15em}
  \begin{alertblock}{Exam tip}
    \footnotesize When asked to explain \emph{why} health care markets fail, trace: characteristic $\rightarrow$ market failure $\rightarrow$ policy instrument. That chain earns full marks.
  \end{alertblock}
\end{frame}

% ----------------------------------------------------------------
\begin{frame}{Review: The Wagstaff Model --- Key Results}
  \begin{block}{Equilibrium condition}
    Individual chooses $H^*$ and $HC^*$ where the indifference curve is tangent to the
    transformed budget constraint in health--consumption space.
  \end{block}
  \vspace{0.4em}
  \begin{center}
    \begin{tabular}{@{}lcc@{}}
      \toprule
      \textbf{Shock} & \textbf{Effect on $HC^*$} & \textbf{Effect on $H^*$} \\
      \midrule
      $\downarrow$ price of HC & $\uparrow$ & $\uparrow$ \\
      $\uparrow$ income (normal good) & $\uparrow$ & $\uparrow$ \\
      Ageing / illness (depreciation $\uparrow$) & $\uparrow$ & $\downarrow$ \\
      $\uparrow$ education (efficiency $\uparrow$) & ambiguous & $\uparrow$ \\
      \bottomrule
    \end{tabular}
  \end{center}
  \vspace{0.3em}
  {\footnotesize Full price of HC = out-of-pocket cost + time and travel costs.
  Insurance lowers the effective price, increasing demand via moral hazard.}
\end{frame}

% ----------------------------------------------------------------
\begin{frame}{Review: Market Failures and Policy Responses}
  \begin{center}
    \begin{tabular}{@{}p{3.5cm}p{4.5cm}p{4cm}@{}}
      \toprule
      \textbf{Market failure} & \textbf{Mechanism} & \textbf{Policy response} \\
      \midrule
      Natural monopoly (insurance admin) & High fixed costs; scale economies & Single public insurer (ACC, Te Whatu Ora) \\[0.3em]
      Moral hazard & Insurance lowers marginal price & Copayments; deductibles; gatekeeping \\[0.3em]
      Adverse selection & Asymmetric info on risk & Community rating; mandates; public provision \\[0.3em]
      Externalities & Spill-overs to others & Subsidise/mandate vaccination; sanitation \\[0.3em]
      SID / asymmetric info & Provider knows more than patient & Licensing; protocols; pay-for-performance \\
      \bottomrule
    \end{tabular}
  \end{center}
\end{frame}

% ----------------------------------------------------------------
\begin{frame}{Review: Health System Design}
  \begin{columns}
    \column{0.32\textwidth}
      \begin{block}{Beveridge model}
        Government funds \emph{and} provides via taxes.\\[0.3em]
        \textbf{NZ, UK}: universal access; monopsony power (PHARMAC); equity goal.
      \end{block}
    \column{0.32\textwidth}
      \begin{alertblock}{Bismarck model}
        Social insurance via payroll; multiple insurers.\\[0.3em]
        \textbf{Germany, Japan, France}: employer--employee co-payment; regulated competition.
      \end{alertblock}
    \column{0.32\textwidth}
      \begin{block}{Private insurance}
        Market-based; fragmented; large admin overhead.\\[0.3em]
        \textbf{US pre-ACA}: adverse selection; uninsured; high spending, poor outcomes.
      \end{block}
  \end{columns}
  \vspace{0.5em}
  \begin{exampleblock}{NZ is a hybrid}
    Beveridge base (Te Whatu Ora, public hospitals) + ACC (no-fault accident insurance)
    + private insurance overlay (Southern Cross, $\approx$33\% of NZ population).
  \end{exampleblock}
\end{frame}

% ----------------------------------------------------------------
\begin{frame}{Practice Question 1 (Short Answer)}
  \begin{columns}
    \column{0.46\textwidth}
      \begin{block}{Question}
        \small Explain the concept of \textbf{supplier-induced demand (SID)} in health care markets.
        What are the theoretical arguments for and against the SID hypothesis, and what
        institutional features of NZ's health system limit or enable SID?
      \end{block}
    \column{0.51\textwidth}
      \begin{alertblock}{Marking guide (3 marks)}
        \small
        \begin{enumerate}
          \item Define SID: providers induce demand beyond patient's ``true'' demand via information asymmetry \hfill[1]
          \item For: fee-for-service; DTC pharma ads; Roemer's Law evidence \hfill[1]
          \item NZ limits: GPs capitation/salaried; PHARMAC; clinical guidelines \hfill[1]
        \end{enumerate}
      \end{alertblock}
  \end{columns}
\end{frame}

% ----------------------------------------------------------------
\begin{frame}{Practice Question 2 (Wagstaff Model)}
  \begin{columns}
    \column{0.43\textwidth}
      \begin{block}{Question}
        \small Use the Wagstaff model to analyse the effect of \textbf{NZ Superannuation eligibility at age 65} on an individual's demand for health care ($HC^*$) and desired health stock ($H^*$). Explain any ambiguities.
      \end{block}
    \column{0.54\textwidth}
      \begin{alertblock}{Marking guide (4 marks)}
        \small
        \begin{enumerate}
          \item Reaching 65: (a) income effect (NZ Super); (b) ageing (depreciation $\uparrow$) \hfill[1]
          \item Income effect: budget shifts out $\Rightarrow$ $HC^*\uparrow$, $H^*\uparrow$ \hfill[1]
          \item Ageing: production fn less efficient $\Rightarrow$ $HC^*\uparrow$, $H^*\downarrow$ \hfill[1]
          \item Net: $HC^*$ almost certainly rises; $H^*$ ambiguous \hfill[1]
        \end{enumerate}
      \end{alertblock}
  \end{columns}
\end{frame}

% ----------------------------------------------------------------
\begin{frame}{Practice Question 3 (Market Failures)}
  \begin{columns}
    \column{0.43\textwidth}
      \begin{block}{Question}
        \small Distinguish \textbf{ex-ante} and \textbf{ex-post} moral hazard. Give a concrete NZ example of each. Explain why moral hazard creates a tension between insurance efficiency and welfare.
      \end{block}
    \column{0.54\textwidth}
      \begin{alertblock}{Marking guide (3 marks)}
        \small
        \begin{enumerate}
          \item Ex-ante: reduced prevention once insured (e.g.\ skipping flu jab; risky behaviour under ACC) \hfill[1]
          \item Ex-post: more HC because OOP price $\approx 0$ (e.g.\ unnecessary GP visits; low-value tests) \hfill[1]
          \item Tension: insurance pools risk, but lower effective price $\Rightarrow$ over-consumption $\Rightarrow$ deadweight loss \hfill[1]
        \end{enumerate}
      \end{alertblock}
  \end{columns}
\end{frame}

% ----------------------------------------------------------------
\begin{frame}{Practice Question 4 (Systems)}
  \begin{columns}
    \column{0.43\textwidth}
      \begin{block}{Question}
        \small Describe the \textbf{US paradox} in health care spending. Use two economic concepts from Part I to explain why the US spends more but achieves worse outcomes than comparable countries.
      \end{block}
    \column{0.54\textwidth}
      \begin{alertblock}{Marking guide (4 marks)}
        \small
        \begin{enumerate}
          \item US paradox: $\approx$17\% GDP on HC --- most of any country --- but last on outcomes (life expectancy, infant mortality) vs.\ OECD \hfill[1]
          \item Admin waste / nat.\ monopoly: fragmented insurance $\Rightarrow$ $\approx$30\% admin overhead vs.\ 12--15\% in single-payer \hfill[1.5]
          \item Adverse selection / uninsured: pre-ACA $\approx$50M uninsured; sicker pool; less preventive care \hfill[1.5]
        \end{enumerate}
      \end{alertblock}
  \end{columns}
\end{frame}

% ----------------------------------------------------------------
\begin{frame}{Common Errors in Part I Exams}
  \begin{enumerate}
    \item \textbf{Confusing moral hazard and adverse selection}\\
          {\small Moral hazard = behaviour changes \emph{after} insurance is purchased.
          Adverse selection = who chooses to buy insurance based on private risk knowledge.}
    \vspace{0.2em}
    \item \textbf{Conflating the six characteristics with market failures}\\
          {\small Characteristics \emph{cause} failures; they are not the same thing.
          Trace the chain: characteristic $\rightarrow$ failure $\rightarrow$ policy.}
    \vspace{0.2em}
    \item \textbf{Forgetting the principal-agent chain in SID}\\
          {\small Always specify: who is the principal, who is the agent, what is the
          information asymmetry, what are the misaligned incentives.}
    \vspace{0.2em}
    \item \textbf{Wagstaff: conflating shifts in the production function with shifts in the
          budget constraint}\\
          {\small Ageing shifts the \emph{production function} (less efficient), not the
          budget constraint. A price change rotates the budget constraint.}
  \end{enumerate}
\end{frame}

% ----------------------------------------------------------------
\begin{frame}{Tomorrow: In-Class Test}
  \begin{alertblock}{Test format and logistics --- Fri 28 Aug 2026}
    \begin{itemize}
      \item \textbf{Weight}: 15\% of final grade
      \item \textbf{Coverage}: all of Part I (Weeks 1--7, including today)
      \item \textbf{Format}: short-answer questions and applied analysis
      \item \textbf{Duration}: full lecture period
      \item No lecture or new reading on Friday --- test-only session
    \end{itemize}
  \end{alertblock}
  \vspace{0.5em}
  \begin{block}{Plussage reminder}
    Grades are computed under (a) the default weights and (b) in-class test down-weighted
    to 0\% with the final exam weighted up correspondingly.
    You receive the higher grade.
  \end{block}
  \vspace{0.3em}
  {\footnotesize After the mid-semester break, Part II begins: Health Economic Evaluation (Week 8, Wed 9 Sep).}
\end{frame}

\summaryside{
  \item Part I synthesised: six characteristics $\rightarrow$ market failures $\rightarrow$
        government involvement; Wagstaff demand model; insurance theory; health system design
  \item Key comparative statics: lower price $\uparrow HC$; ageing $\uparrow HC$ demand
        (depreciation); education $\uparrow H$ (efficiency)
  \item NZ hybrid: Beveridge base + ACC + private insurance overlay; PHARMAC as monopsony purchaser
  \item In-class test tomorrow (Fri 28 Aug, 15\%): no lecture, no new reading; plussage applies
}

\end{document}
